% --- Archon physics formalization source begin ---
% archon:physics
% archon:covers PhyXMiniProblems/problem_phyx_mini_0528.lean
% archon:source-report reports/phyx_mini/problem_phyx_mini_0528.source.json

\chapter{Physics problem phyx\_mini\_0528}
\label{ch:PhyXMiniProblems_problem_phyx_mini_0528}

\paragraph{Problem source.}
\#\# Physical scenario

An observer in a coasting spacecraft moves toward a mirror at speed \textbackslash{}( v = 0.650c \textbackslash{}) relative to the reference frame labeled S in figure. The mirror is stationary with respect to S. A light pulse emitted by the spacecraft travels toward the mirror and is reflected back to the spacecraft. The spacecraft is a distance \textbackslash{}( d = 5.66 \textbackslash{}times 10\textasciicircum{}\{10\} \textbackslash{}) m from the mirror (as measured by observers in S) at the moment the light pulse leaves the spacecraft.

\#\# Question

What is the total travel time of the pulse as measured by observers in the spacecraft?

\#\# Answer choices

- A: A: \textbackslash{}( 77 \textbackslash{}text\{ s\} \textbackslash{})

- B: B: \textbackslash{}( 136 \textbackslash{}text\{ s\} \textbackslash{})

- C: C: \textbackslash{}( 89 \textbackslash{}text\{ s\} \textbackslash{})

- D: D: \textbackslash{}( 174 \textbackslash{}text\{ s\} \textbackslash{})

\#\# Figure caption (auxiliary; use the image as primary evidence)

The image depicts a scene involving a vehicle, a mirror, and labeled elements. 

- **Vehicle**: A futuristic-looking vehicle is shown, moving toward the mirror. An arrow labeled \textbackslash{}(\textbackslash{}vec\{v\}\textbackslash{}) indicates the vehicle's velocity, pointing in the direction of the mirror.

- **Mirror**: The mirror is represented as a vertical blue rectangle on the right side of the image. 

- **Distance**: The distance between the vehicle and the mirror is labeled \textbackslash{}(d\textbackslash{}) with a double-headed arrow indicating measurement.

- **Wave**: A wavy red line with an arrow shows a wave (or signal) moving toward the mirror from the vehicle.

- **Axes**: There is a set of axes with the origin labeled \textbackslash{}(O\textbackslash{}) and the vertical axis labeled \textbackslash{}(S\textbackslash{}).

- **Text**: The words “Mirror,” “\textbackslash{}(d\textbackslash{}),” “\textbackslash{}(\textbackslash{}vec\{v\}\textbackslash{}),” \textbackslash{}(O\textbackslash{}), and \textbackslash{}(S\textbackslash{}) are present, providing labels for the corresponding elements of the diagram.

\#\# Dataset metadata

Relativity; Physical Model Grounding Reasoning, Multi-Formula Reasoning

\paragraph{Recorded answer/context.}
D: D: \textbackslash{}( 174 \textbackslash{}text\{ s\} \textbackslash{})

\paragraph{Figure/image path.}
/root/proposal\_for\_physic/science-mango/phyx\_mini\_run/phyx\_data/test\_image/528.png

\paragraph{Formalization target.}
create a compiling Lean file with sorry bodies at `PhyXMiniProblems/problem\_phyx\_mini\_0528.lean`. The Lean declarations must preserve the physical quantities, dimensions or dimensional roles, figure labels, governing-law hypotheses, and final relation expressed by this problem.
Use Mathlib/Physlib names found through LeanExplore where available. If a physics API is missing, introduce faithful local abstractions rather than scalar placeholder aliases.

\begin{theorem}[Physics formalization target]
\label{thm:physics:phyx_mini_0528:target}
\lean{PhyXMiniProblems.ProblemPhyXMini0528.problem_phyx_mini_0528}
\uses{def:physics:phyx-mini-0528:phyxminiproblems-problemphyxmini0528-spacecraftmirrorroundtripsetup, def:physics:phyx-mini-0528:phyxminiproblems-problemphyxmini0528-matchescoastingspacecraftmirrorscenario, def:physics:phyx-mini-0528:phyxminiproblems-problemphyxmini0528-matchessuppliedspacecraftmirrorfigure, def:physics:phyx-mini-0528:phyxminiproblems-problemphyxmini0528-matchesproblemreadouts, def:physics:phyx-mini-0528:phyxminiproblems-problemphyxmini0528-hasphysicalroundtripparameters, def:physics:phyx-mini-0528:phyxminiproblems-problemphyxmini0528-satisfiesreferenceframelightpropagation, def:physics:phyx-mini-0528:phyxminiproblems-problemphyxmini0528-satisfiesspecialrelativisticpropertime, def:physics:phyx-mini-0528:phyxminiproblems-problemphyxmini0528-recordedanswerchoice, def:physics:phyx-mini-0528:phyxminiproblems-problemphyxmini0528-matchesdisplayedtimechoice, def:physics:phyx-mini-0528:phyxminiproblems-problemphyxmini0528-isuniqueclosesttimechoice}
The assigned autoformalize agent should translate this physics problem into Lean declarations in the covered file, with theorem and lemma proof bodies written as `by sorry`.
\end{theorem}
\begin{proof}
This is an autoformalization task, not a proof task. Produce statements that can later be proved by the physics prover without weakening the physical model.
\end{proof}

% --- Archon declaration topology begin ---
\section{Declaration topology}
\begin{definition}[Length Quantity]
  \label{def:physics:phyx-mini-0528:phyxminiproblems-problemphyxmini0528-lengthquantity}
  \lean{PhyXMiniProblems.ProblemPhyXMini0528.LengthQuantity}
  A nonnegative, unit-independent physical length
\end{definition}
\begin{proof}
  This is the declared model definition of Length Quantity.
\end{proof}
\begin{definition}[Duration Quantity]
  \label{def:physics:phyx-mini-0528:phyxminiproblems-problemphyxmini0528-durationquantity}
  \lean{PhyXMiniProblems.ProblemPhyXMini0528.DurationQuantity}
  A nonnegative, unit-independent physical duration
\end{definition}
\begin{proof}
  This is the declared model definition of Duration Quantity.
\end{proof}
\begin{definition}[Speed Quantity]
  \label{def:physics:phyx-mini-0528:phyxminiproblems-problemphyxmini0528-speedquantity}
  \lean{PhyXMiniProblems.ProblemPhyXMini0528.SpeedQuantity}
  A nonnegative physical speed, used for the spacecraft speed magnitude
\end{definition}
\begin{proof}
  This is the declared model definition of Speed Quantity.
\end{proof}
\begin{definition}[length Readout]
  \label{def:physics:phyx-mini-0528:phyxminiproblems-problemphyxmini0528-lengthreadout}
  \lean{PhyXMiniProblems.ProblemPhyXMini0528.lengthReadout}
  \uses{def:physics:phyx-mini-0528:phyxminiproblems-problemphyxmini0528-lengthquantity}
  Read a physical length in a selected length unit
\end{definition}
\begin{proof}
  This is the declared model definition of length Readout.
\end{proof}
\begin{definition}[duration Readout]
  \label{def:physics:phyx-mini-0528:phyxminiproblems-problemphyxmini0528-durationreadout}
  \lean{PhyXMiniProblems.ProblemPhyXMini0528.durationReadout}
  \uses{def:physics:phyx-mini-0528:phyxminiproblems-problemphyxmini0528-durationquantity}
  Read a physical duration in a selected time unit
\end{definition}
\begin{proof}
  This is the declared model definition of duration Readout.
\end{proof}
\begin{definition}[speed Readout]
  \label{def:physics:phyx-mini-0528:phyxminiproblems-problemphyxmini0528-speedreadout}
  \lean{PhyXMiniProblems.ProblemPhyXMini0528.speedReadout}
  \uses{def:physics:phyx-mini-0528:phyxminiproblems-problemphyxmini0528-speedquantity}
  Read a nonnegative physical speed in selected coherent units
\end{definition}
\begin{proof}
  This is the declared model definition of speed Readout.
\end{proof}
\begin{definition}[vacuum Light Speed Readout]
  \label{def:physics:phyx-mini-0528:phyxminiproblems-problemphyxmini0528-vacuumlightspeedreadout}
  \lean{PhyXMiniProblems.ProblemPhyXMini0528.vacuumLightSpeedReadout}
  Read Physlib's real-valued vacuum speed of light in coherent units
\end{definition}
\begin{proof}
  This is the declared model definition of vacuum Light Speed Readout.
\end{proof}
\begin{definition}[length In Meters]
  \label{def:physics:phyx-mini-0528:phyxminiproblems-problemphyxmini0528-lengthinmeters}
  \lean{PhyXMiniProblems.ProblemPhyXMini0528.lengthInMeters}
  \uses{def:physics:phyx-mini-0528:phyxminiproblems-problemphyxmini0528-lengthquantity, def:physics:phyx-mini-0528:phyxminiproblems-problemphyxmini0528-lengthreadout}
  Metre readout of the initial spacecraft--mirror separation
\end{definition}
\begin{proof}
  This is the declared model definition of length In Meters.
\end{proof}
\begin{definition}[duration In Seconds]
  \label{def:physics:phyx-mini-0528:phyxminiproblems-problemphyxmini0528-durationinseconds}
  \lean{PhyXMiniProblems.ProblemPhyXMini0528.durationInSeconds}
  \uses{def:physics:phyx-mini-0528:phyxminiproblems-problemphyxmini0528-durationquantity, def:physics:phyx-mini-0528:phyxminiproblems-problemphyxmini0528-durationreadout}
  Second readout of a frame-measured physical duration
\end{definition}
\begin{proof}
  This is the declared model definition of duration In Seconds.
\end{proof}
\begin{definition}[speed In Meters Per Second]
  \label{def:physics:phyx-mini-0528:phyxminiproblems-problemphyxmini0528-speedinmeterspersecond}
  \lean{PhyXMiniProblems.ProblemPhyXMini0528.speedInMetersPerSecond}
  \uses{def:physics:phyx-mini-0528:phyxminiproblems-problemphyxmini0528-speedquantity, def:physics:phyx-mini-0528:phyxminiproblems-problemphyxmini0528-speedreadout}
  Metres-per-second readout of the spacecraft's speed magnitude
\end{definition}
\begin{proof}
  This is the declared model definition of speed In Meters Per Second.
\end{proof}
\begin{definition}[vacuum Light Speed In Meters Per Second]
  \label{def:physics:phyx-mini-0528:phyxminiproblems-problemphyxmini0528-vacuumlightspeedinmeterspersecond}
  \lean{PhyXMiniProblems.ProblemPhyXMini0528.vacuumLightSpeedInMetersPerSecond}
  \uses{def:physics:phyx-mini-0528:phyxminiproblems-problemphyxmini0528-vacuumlightspeedreadout}
  Physlib's exact vacuum light speed, read in metres per second
\end{definition}
\begin{proof}
  This is the declared model definition of vacuum Light Speed In Meters Per Second.
\end{proof}
\begin{definition}[Inertial Frame Label]
  \label{def:physics:phyx-mini-0528:phyxminiproblems-problemphyxmini0528-inertialframelabel}
  \lean{PhyXMiniProblems.ProblemPhyXMini0528.InertialFrameLabel}
  The mirror rest frame S and the coasting spacecraft rest frame
\end{definition}
\begin{proof}
  This is the declared model definition of Inertial Frame Label.
\end{proof}
\begin{definition}[Physical Object]
  \label{def:physics:phyx-mini-0528:phyxminiproblems-problemphyxmini0528-physicalobject}
  \lean{PhyXMiniProblems.ProblemPhyXMini0528.PhysicalObject}
  The three physical objects whose worldlines enter the experiment
\end{definition}
\begin{proof}
  This is the declared model definition of Physical Object.
\end{proof}
\begin{definition}[Event Label]
  \label{def:physics:phyx-mini-0528:phyxminiproblems-problemphyxmini0528-eventlabel}
  \lean{PhyXMiniProblems.ProblemPhyXMini0528.EventLabel}
  The three distinguished events of the round trip
\end{definition}
\begin{proof}
  This is the declared model definition of Event Label.
\end{proof}
\begin{definition}[Pulse Leg]
  \label{def:physics:phyx-mini-0528:phyxminiproblems-problemphyxmini0528-pulseleg}
  \lean{PhyXMiniProblems.ProblemPhyXMini0528.PulseLeg}
  The two directed legs of the light pulse
\end{definition}
\begin{proof}
  This is the declared model definition of Pulse Leg.
\end{proof}
\begin{definition}[Axial Direction]
  \label{def:physics:phyx-mini-0528:phyxminiproblems-problemphyxmini0528-axialdirection}
  \lean{PhyXMiniProblems.ProblemPhyXMini0528.AxialDirection}
  Direction along the horizontal spacecraft--mirror axis
\end{definition}
\begin{proof}
  This is the declared model definition of Axial Direction.
\end{proof}
\begin{definition}[Spacecraft Motion Model]
  \label{def:physics:phyx-mini-0528:phyxminiproblems-problemphyxmini0528-spacecraftmotionmodel}
  \lean{PhyXMiniProblems.ProblemPhyXMini0528.SpacecraftMotionModel}
  The spacecraft's stated free inertial motion
\end{definition}
\begin{proof}
  This is the declared model definition of Spacecraft Motion Model.
\end{proof}
\begin{definition}[Mirror Model]
  \label{def:physics:phyx-mini-0528:phyxminiproblems-problemphyxmini0528-mirrormodel}
  \lean{PhyXMiniProblems.ProblemPhyXMini0528.MirrorModel}
  The idealized reflecting object fixed in frame S
\end{definition}
\begin{proof}
  This is the declared model definition of Mirror Model.
\end{proof}
\begin{definition}[Signal Kind]
  \label{def:physics:phyx-mini-0528:phyxminiproblems-problemphyxmini0528-signalkind}
  \lean{PhyXMiniProblems.ProblemPhyXMini0528.SignalKind}
  The electromagnetic signal used by the experiment
\end{definition}
\begin{proof}
  This is the declared model definition of Signal Kind.
\end{proof}
\begin{definition}[Figure Axis]
  \label{def:physics:phyx-mini-0528:phyxminiproblems-problemphyxmini0528-figureaxis}
  \lean{PhyXMiniProblems.ProblemPhyXMini0528.FigureAxis}
  Coordinate lines visible in the supplied image
\end{definition}
\begin{proof}
  This is the declared model definition of Figure Axis.
\end{proof}
\begin{definition}[Figure Label]
  \label{def:physics:phyx-mini-0528:phyxminiproblems-problemphyxmini0528-figurelabel}
  \lean{PhyXMiniProblems.ProblemPhyXMini0528.FigureLabel}
  Semantic transcription of the literal labels visible in the image
\end{definition}
\begin{proof}
  This is the declared model definition of Figure Label.
\end{proof}
\begin{definition}[Figure Arrow]
  \label{def:physics:phyx-mini-0528:phyxminiproblems-problemphyxmini0528-figurearrow}
  \lean{PhyXMiniProblems.ProblemPhyXMini0528.FigureArrow}
  The velocity, pulse, and separation arrows drawn in the image
\end{definition}
\begin{proof}
  This is the declared model definition of Figure Arrow.
\end{proof}
\begin{definition}[Arrow Style]
  \label{def:physics:phyx-mini-0528:phyxminiproblems-problemphyxmini0528-arrowstyle}
  \lean{PhyXMiniProblems.ProblemPhyXMini0528.ArrowStyle}
  Whether a displayed arrow has one or two arrowheads
\end{definition}
\begin{proof}
  This is the declared model definition of Arrow Style.
\end{proof}
\begin{definition}[Figure Color]
  \label{def:physics:phyx-mini-0528:phyxminiproblems-problemphyxmini0528-figurecolor}
  \lean{PhyXMiniProblems.ProblemPhyXMini0528.FigureColor}
  Coarse colors needed to transcribe the salient primary-image evidence
\end{definition}
\begin{proof}
  This is the declared model definition of Figure Color.
\end{proof}
\begin{definition}[Spacecraft Mirror Figure]
  \label{def:physics:phyx-mini-0528:phyxminiproblems-problemphyxmini0528-spacecraftmirrorfigure}
  \lean{PhyXMiniProblems.ProblemPhyXMini0528.SpacecraftMirrorFigure}
  \uses{def:physics:phyx-mini-0528:phyxminiproblems-problemphyxmini0528-physicalobject, def:physics:phyx-mini-0528:phyxminiproblems-problemphyxmini0528-axialdirection, def:physics:phyx-mini-0528:phyxminiproblems-problemphyxmini0528-figureaxis, def:physics:phyx-mini-0528:phyxminiproblems-problemphyxmini0528-figurelabel, def:physics:phyx-mini-0528:phyxminiproblems-problemphyxmini0528-figurearrow, def:physics:phyx-mini-0528:phyxminiproblems-problemphyxmini0528-arrowstyle, def:physics:phyx-mini-0528:phyxminiproblems-problemphyxmini0528-figurecolor}
  Qualitative content of the primary raster. It contains no numerical time or answer-choice value
\end{definition}
\begin{proof}
  This is the declared model definition of Spacecraft Mirror Figure.
\end{proof}
\begin{definition}[Spacecraft Mirror Round Trip Setup]
  \label{def:physics:phyx-mini-0528:phyxminiproblems-problemphyxmini0528-spacecraftmirrorroundtripsetup}
  \lean{PhyXMiniProblems.ProblemPhyXMini0528.SpacecraftMirrorRoundTripSetup}
  \uses{def:physics:phyx-mini-0528:phyxminiproblems-problemphyxmini0528-lengthquantity, def:physics:phyx-mini-0528:phyxminiproblems-problemphyxmini0528-durationquantity, def:physics:phyx-mini-0528:phyxminiproblems-problemphyxmini0528-speedquantity, def:physics:phyx-mini-0528:phyxminiproblems-problemphyxmini0528-inertialframelabel, def:physics:phyx-mini-0528:phyxminiproblems-problemphyxmini0528-physicalobject, def:physics:phyx-mini-0528:phyxminiproblems-problemphyxmini0528-eventlabel, def:physics:phyx-mini-0528:phyxminiproblems-problemphyxmini0528-pulseleg, def:physics:phyx-mini-0528:phyxminiproblems-problemphyxmini0528-axialdirection, def:physics:phyx-mini-0528:phyxminiproblems-problemphyxmini0528-spacecraftmotionmodel, def:physics:phyx-mini-0528:phyxminiproblems-problemphyxmini0528-mirrormodel, def:physics:phyx-mini-0528:phyxminiproblems-problemphyxmini0528-signalkind, def:physics:phyx-mini-0528:phyxminiproblems-problemphyxmini0528-spacecraftmirrorfigure}
  Independent physical quantities, coordinate worldlines, and events of the experiment. The two duration fields are observables in different inertial frames and are not defined from a displayed answer or from the target formula
\end{definition}
\begin{proof}
  This is the declared model definition of Spacecraft Mirror Round Trip Setup.
\end{proof}
\begin{definition}[event Ct Coordinate Meters]
  \label{def:physics:phyx-mini-0528:phyxminiproblems-problemphyxmini0528-eventctcoordinatemeters}
  \lean{PhyXMiniProblems.ProblemPhyXMini0528.eventCtCoordinateMeters}
  \uses{def:physics:phyx-mini-0528:phyxminiproblems-problemphyxmini0528-eventlabel, def:physics:phyx-mini-0528:phyxminiproblems-problemphyxmini0528-spacecraftmirrorroundtripsetup}
  The ct coordinate of an event, interpreted as a real number of metres
\end{definition}
\begin{proof}
  This is the declared model definition of event Ct Coordinate Meters.
\end{proof}
\begin{definition}[event Position In SMeters]
  \label{def:physics:phyx-mini-0528:phyxminiproblems-problemphyxmini0528-eventpositioninsmeters}
  \lean{PhyXMiniProblems.ProblemPhyXMini0528.eventPositionInSMeters}
  \uses{def:physics:phyx-mini-0528:phyxminiproblems-problemphyxmini0528-eventlabel, def:physics:phyx-mini-0528:phyxminiproblems-problemphyxmini0528-spacecraftmirrorroundtripsetup}
  The horizontal x coordinate of an event, in metres in frame S
\end{definition}
\begin{proof}
  This is the declared model definition of event Position In SMeters.
\end{proof}
\begin{definition}[event Time In SSeconds]
  \label{def:physics:phyx-mini-0528:phyxminiproblems-problemphyxmini0528-eventtimeinsseconds}
  \lean{PhyXMiniProblems.ProblemPhyXMini0528.eventTimeInSSeconds}
  \uses{def:physics:phyx-mini-0528:phyxminiproblems-problemphyxmini0528-vacuumlightspeedinmeterspersecond, def:physics:phyx-mini-0528:phyxminiproblems-problemphyxmini0528-eventlabel, def:physics:phyx-mini-0528:phyxminiproblems-problemphyxmini0528-spacecraftmirrorroundtripsetup, def:physics:phyx-mini-0528:phyxminiproblems-problemphyxmini0528-eventctcoordinatemeters}
  Frame-S coordinate time of an event, obtained from ct / c
\end{definition}
\begin{proof}
  This is the declared model definition of event Time In SSeconds.
\end{proof}
\begin{definition}[speed Fraction Of Light]
  \label{def:physics:phyx-mini-0528:phyxminiproblems-problemphyxmini0528-speedfractionoflight}
  \lean{PhyXMiniProblems.ProblemPhyXMini0528.speedFractionOfLight}
  \uses{def:physics:phyx-mini-0528:phyxminiproblems-problemphyxmini0528-speedinmeterspersecond, def:physics:phyx-mini-0528:phyxminiproblems-problemphyxmini0528-vacuumlightspeedinmeterspersecond, def:physics:phyx-mini-0528:phyxminiproblems-problemphyxmini0528-spacecraftmirrorroundtripsetup}
  Dimensionless relative speed β = v/c
\end{definition}
\begin{proof}
  This is the declared model definition of speed Fraction Of Light.
\end{proof}
\begin{definition}[spacecraft Lorentz Factor]
  \label{def:physics:phyx-mini-0528:phyxminiproblems-problemphyxmini0528-spacecraftlorentzfactor}
  \lean{PhyXMiniProblems.ProblemPhyXMini0528.spacecraftLorentzFactor}
  \uses{def:physics:phyx-mini-0528:phyxminiproblems-problemphyxmini0528-spacecraftmirrorroundtripsetup, def:physics:phyx-mini-0528:phyxminiproblems-problemphyxmini0528-speedfractionoflight}
  Physlib's Lorentz factor for the relative spacecraft speed
\end{definition}
\begin{proof}
  This is the declared model definition of spacecraft Lorentz Factor.
\end{proof}
\begin{definition}[emission To Reception Proper Time Seconds]
  \label{def:physics:phyx-mini-0528:phyxminiproblems-problemphyxmini0528-emissiontoreceptionpropertimeseconds}
  \lean{PhyXMiniProblems.ProblemPhyXMini0528.emissionToReceptionProperTimeSeconds}
  \uses{def:physics:phyx-mini-0528:phyxminiproblems-problemphyxmini0528-vacuumlightspeedinmeterspersecond, def:physics:phyx-mini-0528:phyxminiproblems-problemphyxmini0528-spacecraftmirrorroundtripsetup}
  Minkowski proper time between emission and reception, converted from the metre-valued spacetime convention to seconds. This is a general observable definition and contains no problem-specific numerical answer
\end{definition}
\begin{proof}
  This is the declared model definition of emission To Reception Proper Time Seconds.
\end{proof}
\begin{definition}[Matches Coasting Spacecraft Mirror Scenario]
  \label{def:physics:phyx-mini-0528:phyxminiproblems-problemphyxmini0528-matchescoastingspacecraftmirrorscenario}
  \lean{PhyXMiniProblems.ProblemPhyXMini0528.MatchesCoastingSpacecraftMirrorScenario}
  \uses{def:physics:phyx-mini-0528:phyxminiproblems-problemphyxmini0528-spacecraftmirrorroundtripsetup}
  Object roles, frame roles, and directions stated in the prose
\end{definition}
\begin{proof}
  This is the declared model definition of Matches Coasting Spacecraft Mirror Scenario.
\end{proof}
\begin{definition}[Matches Supplied Spacecraft Mirror Figure]
  \label{def:physics:phyx-mini-0528:phyxminiproblems-problemphyxmini0528-matchessuppliedspacecraftmirrorfigure}
  \lean{PhyXMiniProblems.ProblemPhyXMini0528.MatchesSuppliedSpacecraftMirrorFigure}
  \uses{def:physics:phyx-mini-0528:phyxminiproblems-problemphyxmini0528-spacecraftmirrorfigure}
  Primary-image evidence: the spacecraft lies left of the vertical blue mirror; the red velocity and pulse arrows point toward it; d is double-headed; and the labels S, O, Mirror, d, and vector v are visible
\end{definition}
\begin{proof}
  This is the declared model definition of Matches Supplied Spacecraft Mirror Figure.
\end{proof}
\begin{definition}[Matches Problem Readouts]
  \label{def:physics:phyx-mini-0528:phyxminiproblems-problemphyxmini0528-matchesproblemreadouts}
  \lean{PhyXMiniProblems.ProblemPhyXMini0528.MatchesProblemReadouts}
  \uses{def:physics:phyx-mini-0528:phyxminiproblems-problemphyxmini0528-lengthinmeters, def:physics:phyx-mini-0528:phyxminiproblems-problemphyxmini0528-spacecraftmirrorroundtripsetup, def:physics:phyx-mini-0528:phyxminiproblems-problemphyxmini0528-speedfractionoflight}
  The two numerical readouts supplied by the problem statement. They determine the initial separation and relative speed, not either round-trip duration
\end{definition}
\begin{proof}
  This is the declared model definition of Matches Problem Readouts.
\end{proof}
\begin{definition}[Has Physical Round Trip Parameters]
  \label{def:physics:phyx-mini-0528:phyxminiproblems-problemphyxmini0528-hasphysicalroundtripparameters}
  \lean{PhyXMiniProblems.ProblemPhyXMini0528.HasPhysicalRoundTripParameters}
  \uses{def:physics:phyx-mini-0528:phyxminiproblems-problemphyxmini0528-lengthinmeters, def:physics:phyx-mini-0528:phyxminiproblems-problemphyxmini0528-durationinseconds, def:physics:phyx-mini-0528:phyxminiproblems-problemphyxmini0528-spacecraftmirrorroundtripsetup, def:physics:phyx-mini-0528:phyxminiproblems-problemphyxmini0528-eventtimeinsseconds, def:physics:phyx-mini-0528:phyxminiproblems-problemphyxmini0528-speedfractionoflight}
  Positivity, subluminality, and chronological ordering of the experiment
\end{definition}
\begin{proof}
  This is the declared model definition of Has Physical Round Trip Parameters.
\end{proof}
\begin{definition}[Satisfies Reference Frame Light Propagation]
  \label{def:physics:phyx-mini-0528:phyxminiproblems-problemphyxmini0528-satisfiesreferenceframelightpropagation}
  \lean{PhyXMiniProblems.ProblemPhyXMini0528.SatisfiesReferenceFrameLightPropagation}
  \uses{def:physics:phyx-mini-0528:phyxminiproblems-problemphyxmini0528-lengthinmeters, def:physics:phyx-mini-0528:phyxminiproblems-problemphyxmini0528-durationinseconds, def:physics:phyx-mini-0528:phyxminiproblems-problemphyxmini0528-speedinmeterspersecond, def:physics:phyx-mini-0528:phyxminiproblems-problemphyxmini0528-vacuumlightspeedinmeterspersecond, def:physics:phyx-mini-0528:phyxminiproblems-problemphyxmini0528-spacecraftmirrorroundtripsetup, def:physics:phyx-mini-0528:phyxminiproblems-problemphyxmini0528-eventpositioninsmeters, def:physics:phyx-mini-0528:phyxminiproblems-problemphyxmini0528-eventtimeinsseconds}
  Worldlines and coincidences in frame S. The coordinate origin is chosen at emission. The mirror stays at d, the spacecraft coasts at v, the pulse travels at c before and after ideal reflection, and the stored frame-S duration is the emission-to-reception coordinate-time difference
\end{definition}
\begin{proof}
  This is the declared model definition of Satisfies Reference Frame Light Propagation.
\end{proof}
\begin{definition}[Satisfies Special Relativistic Proper Time]
  \label{def:physics:phyx-mini-0528:phyxminiproblems-problemphyxmini0528-satisfiesspecialrelativisticpropertime}
  \lean{PhyXMiniProblems.ProblemPhyXMini0528.SatisfiesSpecialRelativisticProperTime}
  \uses{def:physics:phyx-mini-0528:phyxminiproblems-problemphyxmini0528-durationinseconds, def:physics:phyx-mini-0528:phyxminiproblems-problemphyxmini0528-spacecraftmirrorroundtripsetup, def:physics:phyx-mini-0528:phyxminiproblems-problemphyxmini0528-spacecraftlorentzfactor, def:physics:phyx-mini-0528:phyxminiproblems-problemphyxmini0528-emissiontoreceptionpropertimeseconds}
  The generic proper-time law for two events on an inertially coasting clock. The first field connects the dimensionful spacecraft observable to Physlib's Minkowski interval; the second is the equivalent Lorentz time-dilation law. Neither field supplies the numerical duration requested in this problem
\end{definition}
\begin{proof}
  This is the declared model definition of Satisfies Special Relativistic Proper Time.
\end{proof}
\begin{lemma}[reference Frame Round Trip Duration formula]
  \label{lem:physics:phyx-mini-0528:phyxminiproblems-problemphyxmini0528-referenceframeroundtripduration-formula}
  \lean{PhyXMiniProblems.ProblemPhyXMini0528.referenceFrameRoundTripDuration_formula}
  \uses{def:physics:phyx-mini-0528:phyxminiproblems-problemphyxmini0528-lengthinmeters, def:physics:phyx-mini-0528:phyxminiproblems-problemphyxmini0528-durationinseconds, def:physics:phyx-mini-0528:phyxminiproblems-problemphyxmini0528-speedinmeterspersecond, def:physics:phyx-mini-0528:phyxminiproblems-problemphyxmini0528-vacuumlightspeedinmeterspersecond, def:physics:phyx-mini-0528:phyxminiproblems-problemphyxmini0528-spacecraftmirrorroundtripsetup, def:physics:phyx-mini-0528:phyxminiproblems-problemphyxmini0528-hasphysicalroundtripparameters, def:physics:phyx-mini-0528:phyxminiproblems-problemphyxmini0528-satisfiesreferenceframelightpropagation}
  Solving the two light worldlines and the spacecraft worldline in frame S gives Δt\_S = 2d / (c + v). This is a conclusion, not a propagation-law field
\end{lemma}
\begin{proof}
  The conclusion follows from the governing relations and figure readouts named in the statement of reference Frame Round Trip Duration formula.
\end{proof}
\begin{lemma}[spacecraft Round Trip Duration formula]
  \label{lem:physics:phyx-mini-0528:phyxminiproblems-problemphyxmini0528-spacecraftroundtripduration-formula}
  \lean{PhyXMiniProblems.ProblemPhyXMini0528.spacecraftRoundTripDuration_formula}
  \uses{def:physics:phyx-mini-0528:phyxminiproblems-problemphyxmini0528-lengthinmeters, def:physics:phyx-mini-0528:phyxminiproblems-problemphyxmini0528-durationinseconds, def:physics:phyx-mini-0528:phyxminiproblems-problemphyxmini0528-speedinmeterspersecond, def:physics:phyx-mini-0528:phyxminiproblems-problemphyxmini0528-vacuumlightspeedinmeterspersecond, def:physics:phyx-mini-0528:phyxminiproblems-problemphyxmini0528-spacecraftmirrorroundtripsetup, def:physics:phyx-mini-0528:phyxminiproblems-problemphyxmini0528-spacecraftlorentzfactor, def:physics:phyx-mini-0528:phyxminiproblems-problemphyxmini0528-hasphysicalroundtripparameters, def:physics:phyx-mini-0528:phyxminiproblems-problemphyxmini0528-satisfiesreferenceframelightpropagation, def:physics:phyx-mini-0528:phyxminiproblems-problemphyxmini0528-satisfiesspecialrelativisticpropertime}
  The coasting spacecraft clock records the frame-S round-trip duration divided by γ(v/c). Combining this with the light-worldline result yields the exact proper-time formula used to evaluate the answer list
\end{lemma}
\begin{proof}
  The conclusion follows from the governing relations and figure readouts named in the statement of spacecraft Round Trip Duration formula.
\end{proof}
\begin{definition}[Answer Choice]
  \label{def:physics:phyx-mini-0528:phyxminiproblems-problemphyxmini0528-answerchoice}
  \lean{PhyXMiniProblems.ProblemPhyXMini0528.AnswerChoice}
  Labels of the four whole-second answers in the source problem
\end{definition}
\begin{proof}
  This is the declared model definition of Answer Choice.
\end{proof}
\begin{definition}[seconds]
  \label{def:physics:phyx-mini-0528:phyxminiproblems-problemphyxmini0528-answerchoice-seconds}
  \lean{PhyXMiniProblems.ProblemPhyXMini0528.AnswerChoice.seconds}
  \uses{def:physics:phyx-mini-0528:phyxminiproblems-problemphyxmini0528-answerchoice}
  Seconds printed beside each answer label
\end{definition}
\begin{proof}
  This is the declared model definition of seconds.
\end{proof}
\begin{definition}[recorded Answer Choice]
  \label{def:physics:phyx-mini-0528:phyxminiproblems-problemphyxmini0528-recordedanswerchoice}
  \lean{PhyXMiniProblems.ProblemPhyXMini0528.recordedAnswerChoice}
  \uses{def:physics:phyx-mini-0528:phyxminiproblems-problemphyxmini0528-answerchoice}
  The answer label recorded by the source dataset
\end{definition}
\begin{proof}
  This is the declared model definition of recorded Answer Choice.
\end{proof}
\begin{definition}[Matches Displayed Time Choice]
  \label{def:physics:phyx-mini-0528:phyxminiproblems-problemphyxmini0528-matchesdisplayedtimechoice}
  \lean{PhyXMiniProblems.ProblemPhyXMini0528.MatchesDisplayedTimeChoice}
  \uses{def:physics:phyx-mini-0528:phyxminiproblems-problemphyxmini0528-durationquantity, def:physics:phyx-mini-0528:phyxminiproblems-problemphyxmini0528-durationinseconds, def:physics:phyx-mini-0528:phyxminiproblems-problemphyxmini0528-answerchoice}
  Agreement with a whole-second displayed value. The half-second tolerance is the usual interval for rounding a real duration to the nearest second
\end{definition}
\begin{proof}
  This is the declared model definition of Matches Displayed Time Choice.
\end{proof}
\begin{definition}[Is Unique Closest Time Choice]
  \label{def:physics:phyx-mini-0528:phyxminiproblems-problemphyxmini0528-isuniqueclosesttimechoice}
  \lean{PhyXMiniProblems.ProblemPhyXMini0528.IsUniqueClosestTimeChoice}
  \uses{def:physics:phyx-mini-0528:phyxminiproblems-problemphyxmini0528-durationquantity, def:physics:phyx-mini-0528:phyxminiproblems-problemphyxmini0528-durationinseconds, def:physics:phyx-mini-0528:phyxminiproblems-problemphyxmini0528-answerchoice}
  The chosen value is strictly closer than every alternative in the list
\end{definition}
\begin{proof}
  This is the declared model definition of Is Unique Closest Time Choice.
\end{proof}
% --- Archon declaration topology end ---
% --- Archon physics formalization source end ---
